% SPDX-License-Identifier: Apache-2.0
% covers: Gpio
\datasheet{Gpio}{General purpose pins on AXI-Lite: a value out, a value in, a direction, and an interrupt per pin.}

\dsfacts{\code{lib/parts/src/gpio.rs}}{\code{txhdl\_parts::gpio::Gpio}}%
{\code{//docs:examples}}

\subsection*{Function}

\code{Gpio<N>} puts \code{N} pins behind an AXI-Lite link. A program
chooses which way each pin faces, drives the ones that face out, reads
the ones that face in, and may ask for an interrupt from any of them.

The part drives no pin itself. It offers the three wires a pad wants,
\code{drive}, the value; \code{dirs}, whether to drive it; and
\code{pins}, the value read back, and a board's top level joins them
to a tristate buffer. A \code{Pad} carries nothing in a simulation, so
a part written around one could not be run or checked.

\code{pins} is taken through two flip-flops, \code{sync0} then
\code{sync1}, before anything reads it. Nothing else in the part reads
the port.

\subsection*{Parameters}

\begin{center}\footnotesize
\begin{tabular}{@{}lp{0.7\textwidth}@{}}
\toprule
Parameter & Meaning \\
\midrule
\code{N} & How many pins there are. Every register is \code{N} bits
wide and reads back into the low \code{N} bits of a word; the bits
above are zero. The tables use 8, as \code{ex\_gpio} does. \\
\bottomrule
\end{tabular}
\end{center}

\subsection*{Register map}

Offsets from the peripheral's base. The part selects the word by bits
4 to 2 of the address and looks at no other bits, so the router or the
bridge in front of it decides the base.

\begin{center}\footnotesize
\begin{tabular}{@{}llp{0.5\textwidth}@{}}
\toprule
Offset & Word & Access \\
\midrule
\code{0x00} & \code{out} & Read and write. Driven on a pin whose
direction is out. \\
\code{0x04} & \code{sync1} & Read only. What the pins read, after the
two flip-flops. A write is taken and ignored. \\
\code{0x08} & \code{dir} & Read and write. One to drive the pin, zero
to read it. \\
\code{0x0c} & \code{ie} & Read and write. One to let the pin raise
\code{irq}. \\
\code{0x10} & \code{kind} & Read and write. Zero for a level, one for
an edge. \\
\code{0x14} & \code{pol} & Read and write. A level: one for high. An
edge: one for rising. \\
\code{0x18} & \code{status} & Read, and write one to clear. Which pins
have fired. \\
\bottomrule
\end{tabular}
\end{center}

VHDL reserves \code{out}, so the netlist calls the register
\code{out\_rw} in both targets (issue 497).

\dstables{Gpio}

\subsection*{Behaviour}

\begin{itemize}
\item Every cycle \code{sync0} takes \code{pins}, \code{sync1} takes
\code{sync0}, and \code{seen} takes \code{sync1}. So a read of
\code{sync1} is what the pin held two cycles before, and \code{seen}
is what it held three cycles before, which is what an edge is measured
against.
\item A pin fires while it is enabled in \code{ie} and either its
\code{kind} bit is zero and \code{sync1} equals \code{pol}, or its
\code{kind} bit is one and \code{sync1} crossed \code{seen} in the
direction \code{pol} names.
\item \code{status} takes what it held, less any bit a write to
\code{0x18} set to one in the same cycle, plus every pin firing that
cycle. A level therefore sets its bit again in the cycle after it is
cleared, while the level lasts. Clearing is by a written one and not a
written zero, so two programs clearing different pins cannot lose each
other's work.
\item \code{irq} is high while any bit of \code{status} is set. It is
combinational from \code{status}, which is a register, so it follows
the register by no cycles.
\item A read is answered on \code{bus\_r} in the cycle it is taken, with
\code{OKAY}. A write is taken when its address phase, its beat and
room on \code{bus\_b} are all there, and is answered \code{OKAY} in that
cycle.
\item An address that selects no word reads as zero and a write to it
is answered and does nothing.
\end{itemize}

\subsection*{Verification}

\begin{itemize}
\item \code{//lib/examples:ex\_gpio} sets the direction, drives a
pattern, reads a pin through the synchroniser, takes an interrupt on a
rising edge and not on the fall, clears the status, and asks for a
level and finds the bit back. Every step is asserted against what the
program read.
\item The build simulates \code{Gpio<8>}'s netlist against that run
under nvc and Verilator: \code{//docs:gpio\_sim\_gpio\_tb\_test} and
\code{//docs:gpio\_vsim\_test}.
\end{itemize}

\subsection*{Used by}

Nothing in the tree yet. It is meant for a board's LEDs and keys, and
for the header pins, behind an AXI-Lite bridge; issue 145 says so.

\subsection*{Limits}

The interrupt is one line for all \code{N} pins; which pin raised it
is in \code{status}. There is no debounce, so a key wired to a pin
bounces into an edge interrupt many times. The part does not drive a
pad, so a board that wants one writes the tristate buffer itself.
\code{N} is at most the width of a bus word, since every register
reads back in one.
